% Part: sets-functions-relations
% Chapter: arithmetization
% Section: cuts
%
\documentclass[../../../include/open-logic-section]{subfiles}

\begin{document}

\olfileid{sfr}{arith}{cuts}
\olsection{从 $\Rat$ 到 $\Real$}

本质上，完备性性质表明，实数轴上的任意一点 $\alpha$ 都恰好将该数轴分为两半：那些以 $\alpha$ 为上确界的部分，以及那些以 $\alpha$ 为下确界的部分。为了从有理数 \emph{构造} 实数，Dedekind（戴德金）建议我们直接把实数视为分割有理数的那些\emph{割}。也就是说，我们将 $\sqrt{2}$ 等同于那个把 $< \sqrt{2}$ 的有理数与 $> \sqrt{2}$ 的有理数分开的\emph{割}。

让我们把这一点整理清楚。如果我们将有理数切成两半，那么只需考察其\emph{下半部分}，就能唯一确定该划分。因此，更确切地说，我们给出如下定义：

\begin{defn}[割]
一个\emph{割} $\alpha$ 是指有理数的一个非空、真初始段，并且没有最大元。换言之，$\alpha$ 是一个割当且仅当：
\begin{enumerate}
	\item \emph{非空、真子集}：$\emptyset \neq \alpha \subsetneq \Rat$。
	\item \emph{初始的}：对所有 $p,q \in \Rat$：若 $p < q \in \alpha$ 则 $p \in \alpha$。
	\item \emph{无最大元}：对所有 $p \in \alpha$，存在 $q \in \alpha$ 使得 $p < q$。
\end{enumerate}
那么 $\Real$ 就是所有割的集合。
\end{defn}

现在我们就可以说 $\sqrt{2} = \Setabs{p \in \Rat}{p^2 < 2\text{
or }p < 0}$。当然，我们需要验证这\emph{确实}是一个割，
但我们将此项验证推迟至 \olref[check]{sec}。

如前所述，在定义了某些实体之后，接下来需要定义其上的基本函数与关系。
我们先从一个简单的开始：
\begin{align*}
	\alpha \leq \beta \text{ 当且仅当 }\alpha \subseteq \beta
\end{align*}
这个序的定义使我们能够\emph{陈述}核心结果：割的集合具有完备性性质。
完整展开的话，该命题具有如下形式：若 $S$ 是具有上界的非空割的集合，则 $S$ 有最小上界。
更详细地说：存在一个割 $\lambda$，它是 $S$ 的一个上界，即 $(\forall \alpha \in S)\alpha \subseteq \lambda$，并且 $\lambda$ 是最小的这样的割，即 $(\forall \beta \in \Real)((\forall \alpha \in S)\alpha \subseteq \beta \lif \lambda \subseteq \beta)$。
现在给出这个结果的证明：

\begin{thm}\ollabel{realcompleteness}
割的集合具有完备性性质。
\end{thm}

\begin{proof}
设 $S$ 为任意具有上界的非空割的集合。令 $\lambda = \bigcup S$。
%\Setabs{p \in \Rat}{(\exists \alpha \in S)p \in \alpha}$$
我们首先断言 $\lambda$ 是一个割：
\begin{enumerate}
\item 由于 $S$ 有上界，$S$ 中至少包含一个割，所以 $\emptyset \neq \lambda$。因为 $S$ 是割的集合，所以 $\lambda \subseteq \Rat$。由于 $S$ 有上界，存在某个 $p \in \Rat$ 不在任何一个割 $\alpha \in S$ 中。因此 $p\notin \lambda$，从而 $\lambda \subsetneq \Rat$。
\item 假设 $p < q \in \lambda$。则存在某个 $\alpha \in S$ 使得 $q \in \alpha$。由于 $\alpha$ 是割，$p \in \alpha$。所以 $p \in \lambda$。
\item 假设 $p \in \lambda$。则存在某个 $\alpha \in S$ 使得 $p \in \alpha$。由于 $\alpha$ 是割，存在某个 $q \in \alpha$ 使得 $p < q$。所以 $q \in \lambda$。
\end{enumerate}
这就证明了断言。此外，显然 $(\forall \alpha \in S)\alpha
\subseteq \bigcup S = \lambda$，即 $\lambda$ 是 $S$ 的一个上界。现在假设 $\beta \in \mathbb{R}$ 也是上界，即 $(\forall \alpha \in S)\alpha \subseteq \beta$。对任意 $p \in \Rat$，如果 $p \in \lambda$，则存在 $\alpha \in S$ 使得 $p \in \alpha$，从而 $p \in \beta$。一般地，有 $\lambda \subseteq \beta$。所以 $\lambda$ 是 $S$ 的\emph{最小}上界。
\end{proof}

因此，我们得到了一批满足完备性性质的实体。换种说法就是：我们的割中没有“间隙”。（也就是说，对实数——而非有理数——再作“割”，不会产生任何有趣的新对象。）

接下来，我们必须在实数上定义一些运算。我们首先通过规定 $p_{\Real} =
\Setabs{q \in \Rat}{q < p}$ 对每个 $p \in \Rat$ 成立，将有理数嵌入到实数中。然后定义：
\begin{align*}
	\alpha + \beta &= \Setabs{p + q}{p \in \alpha \land q \in \beta}\\
%	\alpha - \beta &\defis \Setabs{p - q}{p \in \alpha \land q \in \Rat \setminus \beta}\\
	\alpha \times \beta &=
	\Setabs{p \times q}{0 \leq p \in \alpha \land 0 \leq q \in \beta} \cup 0^{\mathbb{R}} & \text{若 }\alpha, \beta \geq 0_{\Real}
%	\alpha \div \beta &\defis
%	\Setabs{\nicefrac{p}{q}}{p \in \alpha \land q \in \Rat \setminus \beta}  & \text{if }\alpha \geq 0_{\Real}, \beta > 0_{\Real}
\end{align*}
为处理其他乘法情形，先令：%that $0_{\Real} \times \alpha = 0_{\Real} = \alpha \times 0_{\Real}$, and add:
\begin{align*}
	-\alpha &= \Setabs{p - q}{p < 0 \land q \notin \alpha}
\end{align*}
然后规定：
\begin{align*}
	\alpha \times \beta &\defis
	\begin{cases}
		\mathord{-}\alpha \times \mathord{-}\beta &\text{若 }\alpha < 0_{\Real}\text{ 且 }\beta < 0_{\Real}\\
		\mathord{-}(\mathord{-}\alpha \times \beta) &\text{若 }\alpha < 0_{\Real} \text{ 且 }\beta > 0_{\Real}\\
		\mathord{-}(\alpha \times \mathord{-}\beta) &\text{若 }\alpha > 0_{\Real} \text{ 且 }\beta < 0_{\Real}
	\end{cases}
\end{align*}
随后我们需要验证这些定义是否总是产生一个割。最后，我们需要进行一番简单（但冗长）的论证，以证明如此定义的割完全符合应有的性质。但我们将所有这些工作推迟至 \olref[check]{sec}。
\end{document}